\documentclass[twocolumn]{aastex631}
\begin{document}
\title{A residual reionization ionizing-photon-budget shortfall for star-forming galaxies at z\~{}8 (beta systematic corner)}
\author{NebulaMind Lab (autonomous pipeline)}
\affiliation{NebulaMind Lab, \url{https://lab.nebulamind.net}}
\begin{abstract}
Reionization ionizing-photon-budget reconciliation at z\~{}8: star-forming galaxies require f\_esc=0.247 (+0.248/-0.126) to close the budget (Madau-Dickinson SFRD, log xi\_ion=25.5+/-0.15, clumping C=2-5, JWST-SFRD tail), versus indirect-proxy-inferred f\_esc=0.050 (+0.076/-0.030) from LzLCS O32/beta calibrations. Median delta(required-inferred)=+0.181 dex-frac (16-84\%: +0.043 to +0.429); 91\% of the systematic MC shows a shortfall. Conclusion: a genuine SHORTFALL remains, and the sign holds under both O32 and beta calibrations. The result is bounded by the xi\_ion x clumping x proxy-calibration systematic, not by statistics. Generated autonomously from public data (jwst) via the NebulaMind Lab runner: a bounded, reproducible, descriptive study.
\end{abstract}
\section{Introduction}
The question of whether star-forming galaxies can provide enough ionizing photons to drive cosmic reionization has been a topic of ongoing debate. Recent studies have highlighted potential shortfalls in the photon budget, suggesting that additional sources or mechanisms may be necessary [Muñoz2024]. This issue is further complicated by uncertainties in key parameters such as the escape fraction (f\_esc) and ionizing efficiency (xi\_ion). To address this problem, we revisit the reionization-photon-budget using a literature-anchored approach.
\section{Data and method}
In our analysis, we adopt the cosmic star formation rate density (SFRD) from Madau \& Dickinson (2014), along with published values for xi\_ion and the O32/beta f\_esc proxy calibrations [LzLCS: Chisholm+22, Flury+22; Simmonds+24]. We do not rely on new observational data or survey catalogs but instead focus on reconciling existing literature values to assess the photon budget at z\~{}8.
\section{Result}
Our calculations indicate that star-forming galaxies require an escape fraction of f\_esc=0.247 (+0.248/-0.126) to close the ionizing-photon-budget, assuming a Madau-Dickinson SFRD and log xi\_ion=25.5±0.15 with clumping factor C=2-5. However, indirect-proxy-inferred values from LzLCS O32/beta calibrations yield f\_esc=0.050 (+0.076/-0.030), resulting in a median shortfall of +0.181 dex-frac (16-84\%: +0.043 to +0.429). Notably, 91\% of our systematic Monte Carlo simulations show a deficit.
\begin{figure}\centering\includegraphics[width=\columnwidth]{result.png}\caption{A residual reionization ionizing-photon-budget shortfall for star-forming galaxies at z\~{}8 (beta systematic corner)}\end{figure}
\section{Caveats}
It is essential to acknowledge the limitations of this approach. Our analysis relies on literature values that may not fully capture the complexities of reionization processes. The use of uncalibrated proxy calibrations and assumptions about clumping factors introduce uncertainties that can affect our results. Furthermore, the automated single-selection method employed here may overlook important variations in galaxy properties and environments. These caveats highlight the need for further research to refine our understanding of cosmic reionization and its underlying mechanisms.
\section*{References}
\footnotesize
[Muoz2024] Muñoz et al. (2024). Reionization after JWST: a photon budget crisis?. bibcode 2024MNRAS.535L..37M \par [Davies2021] Davies et al. (2021). The Predicament of Absorption-dominated Reionization: Increased Demands on Ionizing Sources. bibcode 2021ApJ...918L..35D \par [Park2022] Park et al. (2022). Calibrating excursion set reionization models to approximately conserve ionizing photons. bibcode 2022MNRAS.517..192P \par [Duncan2015] Duncan et al. (2015). Powering reionization: assessing the galaxy ionizing photon budget at z < 10. bibcode 2015MNRAS.451.2030D \par [Madau2017] Madau (2017). Cosmic Reionization after Planck and before JWST: An Analytic Approach. bibcode 2017ApJ...851...50M

\end{document}
